% !TEX root = ../00_Main/Main.tex
\documentclass[../00_Main/Main.tex]{subfiles}
\begin{document}

\section{Comandos de reproducci\'on}\label{ap:repro}

Todos los comandos se ejecutan desde el directorio \texttt{h1/} del
repositorio \texttt{mPES}, con el entorno virtual \texttt{win\_mpes\_env}
activado, Python 3.12, las dependencias de
\texttt{utils/config/requirements.txt} y las variables
\texttt{VIRTUAL\_ENV}, \texttt{PYTHONIOENCODING=utf-8} y
\texttt{TF\_ENABLE\_ONEDNN\_OPTS=0} (Tabla~\ref{tab:repro-commands}). El
entrenamiento final lee los hiperpar\'ametros de
\texttt{inputs/best\_params.json}; la arquitectura del DQN recurrente y
del Transformer, y el n\'umero de episodios de todos los modelos, se
leen de \texttt{config/CONFIG.py}.

\begin{table}[!ht]
\centering
\caption[Comandos de reproducci\'on]{Comandos de optimizaci\'on bayesiana,
entrenamiento final y evaluaci\'on de cada paquete. \texttt{<pkg>} es
cualquiera de los cinco ensambles optimizables.}
\label{tab:repro-commands}
\small
\resizebox{\linewidth}{!}{%
\begin{tabular}{@{}llll@{}}
\toprule
\textbf{Paquete} & \textbf{Optimizaci\'on bayesiana} & \textbf{Entrenamiento final} & \textbf{Evaluaci\'on} \\
\midrule
\texttt{pes\_base} & --- & \texttt{python -m tabular.pes\_base.ext.train\_rl} & \texttt{python -m tabular.pes\_base} \\
\texttt{pes\_ql}   & \texttt{python -m tabular.pes\_ql.ext.optimize\_rl} & \texttt{python -m tabular.pes\_ql.ext.train\_rl} & \texttt{python -m tabular.pes\_ql} \\
\texttt{pes\_dql}  & \texttt{python -m tabular.pes\_dql.ext.optimize\_rl} & \texttt{python -m tabular.pes\_dql.ext.train\_rl} & \texttt{python -m tabular.pes\_dql} \\
\texttt{pes\_dqn}  & \texttt{python -m ml.pes\_dqn.ext.optimize\_dqn} & \texttt{python -m ml.pes\_dqn.ext.train\_dqn} & \texttt{python -m ml.pes\_dqn} \\
\texttt{pes\_rdqn} & \texttt{python -m ml.pes\_rdqn.ext.optimize\_rdqn} & \texttt{python -m ml.pes\_rdqn.ext.train\_rdqn} & \texttt{python -m ml.pes\_rdqn} \\
\texttt{pes\_a2c}  & \texttt{python -m ml.pes\_a2c.ext.optimize\_a2c} & \texttt{python -m ml.pes\_a2c.ext.train\_a2c} & \texttt{python -m ml.pes\_a2c} \\
\texttt{pes\_trf}  & \texttt{python -m ml.pes\_trf.ext.optimize\_tr} & \texttt{python -m ml.pes\_trf.ext.train\_transformer} & \texttt{python -m ml.pes\_trf} \\
\texttt{pes\_ens}  & --- & --- & \texttt{python -m ens.pes\_ens} \\
\texttt{<pkg>}     & \texttt{python -m ens.<pkg>.ext.optimize\_ens} & --- & \texttt{python -m ens.<pkg>} \\
\bottomrule
\end{tabular}
}%
\end{table}

\section{Ejecuci\'on de la evaluaci\'on}\label{ap:orchestrator}

La evaluaci\'on en los $22$ escenarios se ejecuta desde \texttt{h1/}:

\begin{verbatim}
python -m general.scripts.benchmark run --suite both
python -m general.scripts.analysis
python -m general.scripts.figures
python -m general.scripts.random_baseline
\end{verbatim}

El primer comando, reanudable, eval\'ua cada par (modelo, escenario) de
los dos grupos reemplazando temporalmente los archivos de entrada de
cada paquete por los del escenario. El segundo calcula las matrices
modelo $\times$ escenario y las comparaciones entre pares; el tercero
genera las figuras, y el cuarto, las del agente aleatorio. Los resultados
se guardan en \texttt{general/results/<grupo>/}.

Las frecuencias de la Tabla~\ref{tab:ens-freq} se obtienen, desde la
ra\'iz del repositorio, con
\texttt{python writings/auxiliar/scripts/ensemble\_decisions.py}, que
repite la evaluaci\'on del ensamble ponderado, del consenso con
\emph{prior} y de la compuerta del Transformer, registra en cada decisi\'on
la acci\'on antes y despu\'es de cada regla y comprueba que las medias
coincidan con las de la evaluaci\'on principal.

\section{Mapas estad\'isticos complementarios}\label{ap:stat-maps}

Las figuras de este ap\'endice detallan, modelo por modelo, las
comparaciones estad\'isticas que el texto informa de forma agregada.

\begin{figure}[!htbp]
  \centering
  \includegraphics[width=\linewidth]{ind_03_welch_logp_por_escenario.png}
  \caption[Welch por escenario (modelos individuales)]{Modelos
  individuales: $\log_{10} p$ de Welch de cada escenario frente a la
  referencia del mismo modelo; los valores menores que $-10$ se muestran
  como $\le -10$.}
  \label{fig:heatmap-welch}
\end{figure}


\begin{figure}[!htbp]
  \centering
  \includegraphics[width=0.85\linewidth]{ind_13_pares_cohen_d.png}
  \caption[$d$ de Cohen entre pares (modelos individuales)]{Modelos
  individuales: $d$ de Cohen entre pares de modelos en los $21$ escenarios
  de generalizaci\'on.}
  \label{fig:pairwise-cohen}
\end{figure}

\begin{figure}[!htbp]
  \centering
  \begin{subfigure}{\linewidth}
    \centering
    \includegraphics[width=\linewidth]{ens_03_welch_logp_por_escenario.png}
    \caption{$\log_{10} p$ de Welch.}
  \end{subfigure}

  \begin{subfigure}{\linewidth}
    \centering
    \includegraphics[width=\linewidth]{ens_07_cohen_d_por_escenario.png}
    \caption{$d$ de Cohen.}
  \end{subfigure}
  \caption[Mapas estad\'isticos por escenario (ensambles)]{Ensambles:
  comparaci\'on de cada escenario con la referencia del mismo ensamble.}
  \label{fig:ensemble-statistical-heatmaps}
\end{figure}

\begin{figure}[!htbp]
  \centering
  \begin{subfigure}{\linewidth}
    \centering
    \includegraphics[width=0.85\linewidth]{ens_13_pares_cohen_d.png}
    \caption{$d$ de Cohen.}
  \end{subfigure}

  \begin{subfigure}{\linewidth}
    \centering
    \includegraphics[width=0.85\linewidth]{ens_12_pares_welch_logp.png}
    \caption{$\log_{10} p$ de Welch.}
  \end{subfigure}
  \caption[Comparaci\'on entre pares (ensambles)]{Ensambles:
  comparaci\'on entre pares en los $21$ escenarios de generalizaci\'on.}
  \label{fig:ensemble-pairwise}
\end{figure}

\biblio
\end{document}
